\begin{figure}[htbp]
  \centering
  \resizebox{0.85\textwidth}{!}{%
  \begin{tikzpicture}[
    >=Latex,
    font=\small,
    node distance=10mm and 16mm,
    outerbox/.style={
      draw,
      rounded corners=3pt,
      fill=gray!6
    },
    box/.style={
      draw,
      rounded corners=3pt,
      align=center,
      minimum width=30mm,
      minimum height=9mm
    },
    comp/.style={box, fill=blue!6},
    inh/.style={box, fill=green!8},
    arrow/.style={-{Latex[length=2.3mm]}, thick}
  ]
    \node[outerbox, minimum width=56mm, minimum height=54mm] (compgroup) at (-3.8,-0.4) {};
    \node[anchor=north] at (compgroup.north) {\textbf{Composition}};

    \node[comp] (phrase) at (-3.8,1.0) {Phrase};
    \node[comp] (motif) at (-3.8,-0.4) {Motif};
    \node[comp] (note) at (-3.8,-1.8) {NoteEvent};

    \draw[arrow] (phrase) -- node[right] {contains} (motif);
    \draw[arrow] (motif) -- node[right] {contains} (note);

    \node[outerbox, minimum width=76mm, minimum height=54mm] (inhgroup) at (4.2,-0.4) {};
    \node[anchor=north] at (inhgroup.north) {\textbf{Inheritance}};

    \node[inh] (base) at (4.2,0.5) {MusicalObject};
    \node[inh] (melody) at (2.5,-1.2) {Melody};
    \node[inh] (chord) at (5.9,-1.2) {ChordProgression};

    \draw[arrow] (melody) -- (base);
    \draw[arrow] (chord) -- (base);
  \end{tikzpicture}%
  }
  \caption{Comparison of two common relationships in object-oriented design. Composition models part-whole structures such as phrases, motifs, and note events, while inheritance models an is-a relationship where several classes share behaviour from a more general parent class.}
  \label{fig:oop_composition_inheritance}
\end{figure}
